\begin{abstract}
Tensor cores supply most of a GPU's arithmetic throughput, and each NVIDIA generation has bought more of it with silicon: from the V100 to the B200, dense FP16 tensor throughput grew from 125 to 2,200 TFLOPS, while throughput per transistor rose only 1.8$\times$. Whether a tensor-core feature is worth its silicon therefore depends on its microarchitecture, yet NVIDIA documents what each feature delivers, not how, and shipped silicon cannot be reconfigured to measure the effect of a design parameter. This paper presents FourShadow, the white-box counterpart of a commercial tensor core: an open, synthesizable tensor-core microarchitecture on the RISC-V Vortex GPGPU that implements warp-group MMA and 2:4 structured sparsity with internals that can be read and changed. Where FourShadow agrees with an RTX 5080 and an H100, its mechanism explains theirs: a warp-group MMA waits for its shared-memory $A$ operand on every instruction, a fixed delay independent of the tile width, and 2:4 sparsity halves only part of a GEMM's per-$K$ cost, a part that shrinks as the tile widens. Sweeping FourShadow shows that per-lane throughput peaks at 16 threads per warp and that a wider warpgroup speeds up the warp-group MMA but not the rest of the kernel, which the core's instruction front end limits. On an AMD Alveo U55C FPGA at 250\,MHz, warp-group MMA raises the end-to-end speedup over the SIMT baseline to 9.15$\times$ for Llama~2 prefill and 2.00$\times$ for decode.
\end{abstract}
